\subsection{极大子模}

定义 2. —— 设 A 是环，M 是一个 A-模。M 的极大子模就是 M 的真子模组成的集合（按包含关系排序）中的一个极大元。

\(A_{s}\) 的极大子模只不过是 A 的极大左理想。

设 N 是 M 的一个子模。M/N 的子模都具有形式 P/N，其中 P 是 M 的包含 N 的子模（I, §4, No.~6, 第 41 页，定理 4）。因此，N 是 M 的极大子模当且仅当模 M/N 是单模。

命题 3. —— 设 M 是一个有限生成 A-模。M 的每个真子模都包含在某个极大子模中。

设 N 是 M 的一个真子模。用 S 表示 M 的包含 N 的真子模组成的集合，按包含关系排序；我们来证明 S 是归纳的。设 F 是 S 的一个全序子集。若 F 为空，则 N 是 F 在 S 中的上界。在相反的情形，用 Q 表示 F 的诸元素之并。则 Q 是 M 的一个子模。设 F 是 M 的一个有限生成子集。若 Q 等于 M，则 F 将包含于某个子模 \(P \in F\)，这将蕴含 P = M，与 F 的定义矛盾。于是我们有 \(Q \in S\)，这证明 S 是归纳的。命题 3 于是由《集合论》，III, §2, No.~4, 第 155 页的推论 1 得出。

推论 1. —— 设 M 是一个非零的有限生成 A-模。存在 A 的一个双边理想 a——某个单 A-模的零化子——使得 aM 不等于 M。

设 N 是 M 的一个极大子模（命题 3），a 是单 A-模 M/N 的零化子；它是 A 的一个双边理想，且我们有 \(\mathfrak{a}(M/N)=0\)，由此 \(aM\subset N\)，从而 \(aM\neq M\) 。

推论 2. —— 设 A 是交换环，B 是一个 A-代数。设 M 是一个作为 A-模为有限生成的单 B-模，m 是 A-模 M 的零化子。则 m 是 A 的极大理想，且 M 是体 A/m 上的有限维向量空间。

由于环 A 交换，单 A-模的零化子是 A 的极大理想（VIII，第 46 页，命题 1）。把推论 1 应用于 A-模 M，存在 A 的一个极大理想 a，使 \(aM \neq M\) 。但 aM 是单 B-模 M 的一个子模；于是我们有 aM = 0，从而 \(a \subset m\) 。由于理想 m 异于 A，它等于 a，故 m 是极大的。模 M 是体 A/m 上的有限生成模；最后的断言由此得出。

